% !TeX root = ../navier-stokes-manuscript.tex
% ============================================================
% 2.1 Vortex Stretching
% ============================================================

Axial stretching lengthens one narrow material vortex-tube segment while its material volume stays fixed, so its cross-section contracts, its rotating fluid moves inward, and its interior spins faster. The velocity field carrying the segment obeys the momentum equation
\[
  \underbrace{\partial_t u}_{\text{local acceleration}}
  +
  \underbrace{(u\cdot\nabla)u}_{\text{nonlinear transport}}
  =
  \underbrace{-\nabla p}_{\text{pressure}}
  +
  \underbrace{\nu\Delta u}_{\text{viscous diffusion}}.
\]
Taking the curl removes pressure and turns nonlinear transport into vorticity advection and stretching.

\subsection{Vortex Stretching}\label{sub:ThePerfectlyStretchingVortex}

Place the segment in a local axisymmetric aligned case, with strain locally uniform across the material segment and aligned with its axis. If \(e\) is that axis and \(L(t)\) is the segment's length, write
\[
  S:=\frac12\bigl(\nabla u+(\nabla u)^{\mathsf T}\bigr),
  \qquad
  Se=\sigma(t)e,
  \qquad
  \sigma(t)>0,
  \qquad
  \frac{\dot L(t)}{L(t)}
  =
  e\cdot Se
  =
  \sigma(t).
\]
Incompressibility keeps its material volume \(\mathcal V\) fixed. With effective transverse area \(A(t):=\mathcal V/L(t)\) and effective radius \(r_{\mathrm{eff}}(t):=\sqrt{A(t)/\pi}\),
\[
  \nabla\cdot u=0,
  \qquad
  \frac{d}{dt}(A L)=0,
  \qquad
  \frac{\dot A}{A}=-\sigma(t),
  \qquad
  \frac{\dot r_{\mathrm{eff}}}{r_{\mathrm{eff}}}
  =
  -\frac{\sigma(t)}{2}.
\]
The same narrowing segment carries vorticity \(\omega:=\nabla\times u\), governed by
\[
  \frac{D\omega}{Dt}
  =
  S\omega+\nu\Delta\omega,
  \qquad
  \frac{D}{Dt}
  :=
  \partial_t+u\cdot\nabla,
\]
and under the alignment \(\omega=|\omega|e\),
\[
  S\omega
  =
  \sigma(t)\omega,
  \qquad
  \frac12\frac{D|\omega|^2}{Dt}
  =
  \sigma(t)|\omega|^2
  +
  \nu\,\omega\cdot\Delta\omega.
\]
Aligned stretching drives the rotating fluid toward spin-up, but the segment's squared vorticity grows only where this positive contribution and the signed viscous contribution sum to a positive number. The viscous contribution has the sign of \(\omega\cdot\Delta\omega\). The energy identity proved in \Cref{prop:ClassicalKineticEnergyIdentity} and the antisymmetric part of \(\nabla u\) give
\[
  \norm{u(t)}_2
  \leq
  \norm{u_0}_2
  <\infty,
  \qquad
  \left|\frac12\bigl(\nabla u-(\nabla u)^{\mathsf T}\bigr)\right|_{\mathrm F}
  =
  \frac{|\omega|}{\sqrt2}
  \leq
  |\nabla u|_{\mathrm F}.
\]
Finite kinetic energy can coexist with concentration in the rotational part of the velocity gradient inside the narrowing segment, so the full field can be compared directly at finer width.

\divider{\NavierStokes{} Scaling}

At a time \(t_0<T_*\), call the effective radius \(\ell:=r_{\mathrm{eff}}(t_0)\). For \(\lambda>1\), define
\[
  u_\lambda(x,t)
  :=
  \lambda u(\lambda x,\lambda^2t),
  \qquad
  p_\lambda(x,t)
  :=
  \lambda^2p(\lambda x,\lambda^2t).
\]
On the corresponding rescaled interval, \(u_\lambda\) is another \NavierStokes{} solution. At the corresponding rescaled time \(\lambda^{-2}t_0\), its segment has width \(\lambda^{-1}\ell\). This rescaled field is not a later state of the opening material segment. Its momentum terms scale as
\[
\begin{aligned}
  \partial_tu_\lambda(x,t)
  &=
  \lambda^3(\partial_tu)(\lambda x,\lambda^2t),
  \\
  (u_\lambda\cdot\nabla)u_\lambda(x,t)
  &=
  \lambda^3\bigl((u\cdot\nabla)u\bigr)(\lambda x,\lambda^2t),
  \\
  \nabla p_\lambda(x,t)
  &=
  \lambda^3(\nabla p)(\lambda x,\lambda^2t),
  \\
  \nu\Delta u_\lambda(x,t)
  &=
  \lambda^3\nu(\Delta u)(\lambda x,\lambda^2t).
\end{aligned}
\]
Rescaling multiplies every momentum term by the same cubic factor, so neither nonlinear transport nor viscosity gains relative ground.

\divider{Critical Height}

The Sobolev norms of the full field change with order. Let \(\Lambda:=(-\Delta)^{1/2}\). Since
\[
  \widehat{u_\lambda}(\xi,t)
  =
  \lambda^{-2}\widehat u(\lambda^{-1}\xi,\lambda^2t),
\]
a change of variables gives
\[
  \norm{u_\lambda(t)}_{\dot H^s}^2
  =
  \lambda^{2s-1}
  \norm{u(\lambda^2t)}_{\dot H^s}^2.
\]
At \(s=0\), kinetic energy loses one power of \(\lambda\). At \(s=1\), the first derivatives gain one power. At \(s=\tfrac12\), the power is zero. This scale-neutral middle defines the critical height
\[
  H(t)
  :=
  \frac12\norm{u(t)}_{\dot H^{1/2}}^2
  =
  \frac12\norm{\Lambda^{1/2}u(t)}_2^2.
\]
Writing \(H_\lambda\) for the same quantity evaluated on \(u_\lambda\),
\[
  \norm{u_\lambda(t)}_2^2
  =
  \lambda^{-1}\norm{u(\lambda^2t)}_2^2,
  \qquad
  H_\lambda(t)=H(\lambda^2t),
  \qquad
  \norm{\nabla u_\lambda(t)}_2^2
  =
  \lambda\norm{\nabla u(\lambda^2t)}_2^2.
\]
Scale neutrality does not determine whether \(H\) rises or falls in time.

\divider{Nonlinear Production versus Viscous Dissipation}

At critical height, nonlinear transport contributes the signed production
\[
  \mathcal P_{1/2}[u](t)
  :=
  -\operatorname{Re}
  \pair
    {\Lambda^{1/2}u(t)}
    {\Lambda^{1/2}\bigl((u(t)\cdot\nabla)u(t)\bigr)}_{L^2}.
\]
Pressure vanishes from the pairing by incompressibility, while viscosity contributes the negative dissipation term \(-\nu\norm{\Lambda^{3/2}u}_2^2\). Thus
\[
  H'(t)
  +
  \nu\norm{\Lambda^{3/2}u(t)}_2^2
  =
  \mathcal P_{1/2}[u](t).
\]
The height rises when production exceeds dissipation and falls when dissipation exceeds production. Under rescaling,
\[
  \mathcal P_{1/2}[u_\lambda](t)
  =
  \lambda^2\mathcal P_{1/2}[u](\lambda^2t),
  \qquad
  \nu\norm{\Lambda^{3/2}u_\lambda(t)}_2^2
  =
  \lambda^2\nu\norm{\Lambda^{3/2}u(\lambda^2t)}_2^2.
\]
Rescaling multiplies production and dissipation by the same quadratic factor. Neither gains ground, and the duration of a nonlinear transition remains undetermined.

\divider{The Heat Clock}

The heat equation \(v_t=\nu\Delta v\) gives a linear viscous comparison:
\[
  \widehat v(\xi,t+\delta t)
  =
  e^{-\nu|\xi|^2\delta t}\widehat v(\xi,t).
\]
A scale-localized velocity structure of width \(r\) occupies Fourier frequencies with magnitude comparable to \(r^{-1}\). Its viscous change becomes order one when
\[
  \nu|\xi|^2\delta t\simeq1,
\]
which gives the linear heat comparison time
\[
  \tau_\nu(r)
  :=
  \frac{r^2}{\nu}.
\]
At a width smaller by \(\lambda^{-1}\), this time becomes
\[
  \tau_\nu(\lambda^{-1}r)
  =
  \lambda^{-2}\tau_\nu(r).
\]
The linear comparison time contracts in the same proportion as time under the full \NavierStokes{} rescaling.

\Needspace{18\baselineskip}
\divider{The Zeno Cascade}

The dyadic widths and their heat times are
\[
  r_n:=2^{-n}r_0,
  \qquad
  \tau_n:=\frac{r_n^2}{\nu},
\]
The heat times shrink geometrically:
\[
  \tau_n
  =
  \frac{r_0^2}{\nu}4^{-n},
  \qquad
  \sum_{n=0}^{\infty}\tau_n
  =
  \frac{4r_0^2}{3\nu}
  <
  \infty.
\]
This finite sum can fit inside a finite interval, but arithmetic does not realize a history. One \NavierStokes{} field would have to exhibit scale-localized structures at the widths
\[
  r_0>r_1>r_2>\cdots\longrightarrow0
\]
and sampled times
\[
  t_0<t_1<t_2<\cdots\uparrow T_*<\infty,
  \qquad
  t_{n+1}-t_n\asymp\tau_n,
\]
with comparison constants independent of \(n\). That same field passes through successively smaller structures in gaps comparable to \(r_n^2/\nu\), while its remaining time satisfies
\[
  T_*-t_n
  \asymp
  \sum_{k=n}^{\infty}\tau_k
  =
  \frac{4r_n^2}{3\nu}.
\]
Those gaps collapse at the finite endpoint. Its first and successive Eulerian time derivatives at one fixed spatial point remain uncontrolled.
